\documentclass[11pt]{article}
\usepackage[margin=1in]{geometry}
\usepackage{array}
\usepackage{booktabs}
\usepackage{longtable}
\usepackage[table]{xcolor}
\usepackage{hyperref}

\title{\texttt{kiss}: Structural Feedback for Agent-Written Code}
\author{GPT-5.5, edited by David Sweet\thanks{\texttt{dsweet@andamooka.org}}}
\date{}

\begin{document}
\maketitle

\section{Design philosophy}

\texttt{kiss} treats code quality feedback as a set of fast, named, falsifiable
signals rather than as a single quality score.  This follows Fowler's account of
code smells: a smell is a quick surface indication that often points to a deeper
problem, but smells are not inherently bad and require inspection
\cite{fowler-code-smell}.  The tool is therefore best understood as a guardrail
for review and refactoring, not as an oracle.

The non-rejected design hypothesis is that local edits can degrade global
structure unless an editor receives compact repository-level feedback.  This is
consistent with SWE-bench, where real GitHub issues often require coordinated
changes across functions, classes, and files \cite{jimenez-swebench}, and with
CodeAgent, where repository tools for retrieval, symbol navigation, and testing
improved repository-level code generation \cite{zhang-codeagent}.  The hypothesis
must be narrowed, because repository context is not automatically helpful:
Gloaguen et al. found that AGENTS.md-style context files can increase inference
cost and, in some settings, reduce success, so feedback should be justified,
targeted, and minimal
\cite{gloaguen-agents}.

The second non-rejected hypothesis is that internal design quality lowers the
cost of future change.  Fowler argues that internal quality matters because
developers spend most of their time modifying existing code, and poor internal
quality increases both understanding time and mistake risk
\cite{fowler-quality}.  Parnas gives the classic modularity criterion:
decompose systems around hidden design decisions so that changes remain local
\cite{parnas}.  Ousterhout similarly frames software design as complexity
management through deep modules and small interfaces \cite{ousterhout}.

\section{Non-rejected and rejected hypotheses}

\begin{longtable}{p{0.43\linewidth}p{0.43\linewidth}}
\toprule
Hypothesis & Status \\
\midrule
\endhead
\rowcolor{gray!18}
Global structural feedback helps agents & Non-rejected, narrowed \\
\multicolumn{2}{p{0.86\linewidth}}{
Use compact, task-relevant feedback.  Repository-scale tasks need context, but
unfiltered context can add cost or hurt
\cite{jimenez-swebench,zhang-codeagent,gloaguen-agents}.} \\[0.35em]

\rowcolor{gray!18}
Size and complexity matter & Non-rejected, calibrated \\
\multicolumn{2}{p{0.86\linewidth}}{
Large or intricate units are useful review targets.  Antinyan et al. report that
some code-complexity characteristics have a major influence on maintenance time,
while Sjoberg et al. found size and low cohesion to be the strongest observed
maintenance-effort signals in their case study \cite{antinyan,sjoberg}.} \\[0.35em]

\rowcolor{gray!18}
Dependency structure matters & Non-rejected, calibrated \\
\multicolumn{2}{p{0.86\linewidth}}{
Fan-out, cycles, and deep dependency chains are risk signals.  Coupling and
structural dependencies have evidence as defect or change-propagation signals,
but that evidence is strongest for coupling-related measures, and fan-in can
also indicate a useful stable abstraction
\cite{elemam,gerosa,parnas}.} \\[0.35em]

\rowcolor{gray!18}
Duplication is always bad & Rejected as an absolute \\
\multicolumn{2}{p{0.86\linewidth}}{
Duplication is a smell, not proof of bad code.  Beck's simple design and the DRY
principle discourage duplicated knowledge \cite{fowler-beck,pragmatic-dry}, but
clone studies show many clones are not more defect-prone \cite{rahman-clones}.} \\[0.35em]

\rowcolor{gray!18}
Test-reference coverage proves quality & Rejected as a quality target \\
\multicolumn{2}{p{0.86\linewidth}}{
Tests are central to safe change \cite{fowler-refactoring,feathers}, but coverage
is only weak-to-moderately related to test effectiveness and should not be used
as a quality target \cite{inozemtseva,kochhar}.  Static name references are an
even coarser signal.} \\[0.35em]

\rowcolor{gray!18}
Semantic refactoring is safer than broad text replacement & Non-rejected \\
\multicolumn{2}{p{0.86\linewidth}}{
Fowler defines refactoring as small behavior-preserving transformations, and
Feathers recommends tests as the safety net for changing legacy code
\cite{fowler-refactoring,feathers}.} \\[0.35em]

\rowcolor{gray!18}
Fast, prioritized feedback matters & Non-rejected, partly extrapolated \\
\multicolumn{2}{p{0.86\linewidth}}{
Fast feedback is consistent with small-step refactoring, and CI evidence shows
that long build durations are a practical problem \cite{fowler-refactoring,ghaleb}.
Static-analysis studies also show warnings must be prioritized and useful, not
merely numerous \cite{vassallo,motwani}.} \\[0.35em]

\rowcolor{gray!18}
One universal code-health score is enough & Rejected \\
\multicolumn{2}{p{0.86\linewidth}}{
Fowler's smells and Beck's rules are component signals requiring judgment, not a
single score \cite{fowler-code-smell,fowler-beck}.  Maintainability metrics can
be mutually inconsistent, so decomposed signals are easier to inspect and act on
\cite{sjoberg}.} \\[0.35em]
\bottomrule
\end{longtable}

\section{Metrics and confidence}

The confidence column means confidence that the metric is useful as a review or
refactoring signal.  It does not mean that violating the metric proves the code
is bad.  Code smells require contextual judgment \cite{fowler-code-smell}, and
static-analysis warnings are more useful when they are prioritized for
actionability \cite{vassallo,motwani}.

\begin{itemize}
\item \textbf{High confidence} The metric is a strong review signal. There is
direct support from software-engineering research or well-established
practice that the measured structure often affects maintainability, change risk,
or comprehension.
\item \textbf{Medium confidence} The metric is a useful heuristic. The
underlying concern is well supported, but the specific measurement is an
indirect proxy, context-sensitive, or easy to misuse if treated mechanically.
\item \textbf{Low confidence} The metric is a weak or experimental signal. It
may catch real problems, but the evidence is indirect, the proxy is noisy, or
the false-positive risk is high.
\end{itemize}

\begin{longtable}{p{0.43\linewidth}p{0.43\linewidth}}
\toprule
Area & Confidence \\
\midrule
\endhead
\rowcolor{gray!18}
Function size and control-flow complexity &
High \\
\multicolumn{2}{@{\hspace{1.25em}}p{0.86\linewidth}}{\texttt{statements\_per\_function},
\texttt{branches\_per\_function}, \texttt{max\_indentation},
\texttt{nested\_function\_depth}, \texttt{calls\_per\_function}} \\
\multicolumn{2}{p{0.86\linewidth}}{Flags functions that may
require too much local reasoning.  Complexity is a
central construction concern, and empirical work links some complexity
characteristics to maintenance time
\cite{antinyan,mcconnell,fowler-code-smell}.} \\[0.35em]

\rowcolor{gray!18}
Function interface complexity &
Medium \\
\multicolumn{2}{@{\hspace{1.25em}}p{0.86\linewidth}}{\texttt{positional\_args},
\texttt{keyword\_only\_args}, \texttt{arguments},
\texttt{boolean\_parameters}} \\
\multicolumn{2}{p{0.86\linewidth}}{Flags call sites whose
meaning may be hard to read or easy to misuse.  This
aligns with design guidance that code should reveal intention and reduce reader
burden \cite{fowler-beck,ousterhout}.} \\[0.35em]

\rowcolor{gray!18}
Function exit and data-shape complexity &
Medium \\
\multicolumn{2}{@{\hspace{1.25em}}p{0.86\linewidth}}{\texttt{returns\_per\_function},
\texttt{local\_variables}, \texttt{return\_values\_per\_function}} \\
\multicolumn{2}{p{0.86\linewidth}}{Flags functions with many
paths or many local concepts.  These are conservative
review prompts grounded in construction guidance and broader complexity evidence
\cite{antinyan,mcconnell}.} \\[0.35em]

\rowcolor{gray!18}
Exception or annotation surface &
Low--Medium \\
\multicolumn{2}{@{\hspace{1.25em}}p{0.86\linewidth}}{\texttt{statements\_per\_try\_block},
\texttt{decorators\_per\_function}, \texttt{attributes\_per\_function}} \\
\multicolumn{2}{p{0.86\linewidth}}{Keeps exceptional or
meta-programmed regions narrow enough to inspect.  This is
mainly a smell heuristic: broad wrappers and hidden behavior need review, but
the metric is less directly validated than size or coupling \cite{fowler-code-smell}.} \\[0.35em]

\rowcolor{gray!18}
Type or class size &
Medium \\
\multicolumn{2}{@{\hspace{1.25em}}p{0.86\linewidth}}{\texttt{methods\_per\_class},
\texttt{interface\_types\_per\_file}, \texttt{concrete\_types\_per\_file}} \\
\multicolumn{2}{p{0.86\linewidth}}{Flags large type surfaces
and crowded files.  It can reveal god-object or
over-abstraction smells, but deep modules may legitimately contain substantial
implementation behind a small interface \cite{ousterhout,fowler-code-smell}.} \\[0.35em]

\rowcolor{gray!18}
File size and file crowding &
Medium \\
\multicolumn{2}{@{\hspace{1.25em}}p{0.86\linewidth}}{\texttt{statements\_per\_file},
\texttt{lines\_per\_file}, \texttt{functions\_per\_file}} \\
\multicolumn{2}{p{0.86\linewidth}}{Helps find files that have
become change magnets or mixed-concern modules.
System size and local size are useful signals, but splitting files blindly can
increase indirection \cite{sjoberg,ousterhout}.} \\[0.35em]

\rowcolor{gray!18}
Import and dependency graph &
High for cycles/fan-out; \newline Medium for counts \\
\multicolumn{2}{@{\hspace{1.25em}}p{0.86\linewidth}}{\texttt{imported\_names\_per\_file},
\texttt{cycle\_size}, \texttt{indirect\_dependencies},
\texttt{dependency\_depth}} \\
\multicolumn{2}{p{0.86\linewidth}}{Highlights coupling,
cycles, and long chains as review prompts.  Information
hiding is central to modular design, while dependency evidence is strongest for
coupling and change propagation \cite{parnas,elemam,gerosa}.} \\[0.35em]

\rowcolor{gray!18}
Duplication gate &
Medium \\
\multicolumn{2}{@{\hspace{1.25em}}p{0.86\linewidth}}{\texttt{min\_similarity},
\texttt{duplication\_enabled}} \\
\multicolumn{2}{p{0.86\linewidth}}{Finds similar chunks that
may encode duplicated knowledge.  DRY and Beck's
simple design support removing duplicated knowledge \cite{fowler-beck,pragmatic-dry},
but clone evidence says not all clones are harmful \cite{rahman-clones}.} \\[0.35em]

\rowcolor{gray!18}
Static test-name reference gate &
Low as quality proof; \newline Medium as guardrail \\
\multicolumn{2}{@{\hspace{1.25em}}p{0.86\linewidth}}{\texttt{test\_coverage\_threshold}} \\
\multicolumn{2}{p{0.86\linewidth}}{Flags code units whose
names do not appear in tests.  This is not runtime
coverage and should not be read as evidence that a unit is tested.  It is a cheap
pre-refactoring proxy: an unmentioned name suggests the safety net is not
visible, while a mentioned name may still be weakly asserted or incidental.
Runtime coverage studies support only the narrower lesson that coverage can help
find under-tested areas, while warning that coverage is not a reliable quality
target \cite{inozemtseva,kochhar}.  This static proxy is weaker and needs direct
validation on agent edits.} \\[0.35em]

\rowcolor{gray!18}
Orphan module gate &
Low--Medium \\
\multicolumn{2}{@{\hspace{1.25em}}p{0.86\linewidth}}{\texttt{orphan\_module\_enabled}} \\
\multicolumn{2}{p{0.86\linewidth}}{Finds modules that appear
disconnected from the dependency graph.  Orphans may
be dead code or intentional entry points; inspect before deleting.  The value is
rooted in modularity and changeability rather than direct empirical validation
\cite{parnas,fowler-quality}.} \\[0.35em]
\bottomrule
\end{longtable}

\texttt{inv\_test\_coverage} is the inverted static test-reference metric behind
the static test-name reference gate.  In \texttt{kiss stats}, each value is one
source file's percentage of definitions whose names do not appear in tests.  The
following columns are the same distribution columns printed by
\texttt{kiss stats}: \texttt{N} is the number of measured files, \texttt{p50}
through \texttt{p99} are distribution percentiles, and \texttt{max} is the
largest observed value.  Because the metric is inverted, this table discusses
\emph{median static test-reference coverage}: the estimated percentage of
definitions in the median file whose names do appear in tests.  Numerically,
that is \texttt{100 - p50}.

\begingroup
\small
\setlength{\tabcolsep}{3pt}
\begin{longtable}{p{0.12\linewidth}r r r r r r p{0.14\linewidth}r r}
\toprule
Project & N & p50 & p90 & p95 & p99 & max & Coverage & \texttt{kiss} time &
Cov. time \\
\midrule
\endhead
\texttt{kiss} & 302 & 0 & 0 & 9 & 20 & 75 & 92.60\% \texttt{llvm-cov} &
0.3s & 1.5s \\
\texttt{ripgrep} & 100 & 50 & 100 & 100 & 100 & 100 & 60.47\% \texttt{llvm-cov} &
0.4s & 15s \\
\texttt{rich} & 213 & 0 & 100 & 100 & 100 & 100 & 91\% \texttt{slipcover} &
0.2s & 11s \\
\texttt{ruff} & 2917 & 0 & 0 & 100 & 100 & 100 & 89.48\% \texttt{llvm-cov} &
3.5s & 24m \\
\texttt{ruff} & 1828 & 100 & 100 & 100 & 100 & 100 & 62\% \texttt{slipcover} &
-- & 3s \\
\bottomrule
\end{longtable}
\endgroup

These measurements are not a validation study, but they show the intended
tradeoff.  \texttt{kiss} is considerably faster than runtime line-level coverage
measurement, especially when the coverage run requires full project test
execution.  In this small sample, median static test-reference coverage is often
in the same rough range as the line-level coverage number, but the second
\texttt{ruff} row is a counterexample and the metric remains only a static
name-reference proxy.  The \texttt{kiss} repository is a special case: it was
fully agent-coded and has been maintained with \texttt{kiss} as a linter.  Its
high line-level coverage is evidence that agent-written code developed with
\texttt{kiss} feedback can maintain a strong runtime test suite, but not
independent evidence that the static proxy will generalize across projects.

\section{Limitations}

Much of the evidence behind these metrics comes from human software development:
human estimates of maintenance time, human maintainability case studies, human
static-analysis workflows, and human refactoring practice
\cite{antinyan,sjoberg,vassallo,motwani,fowler-refactoring}.  \texttt{kiss}
extrapolates from that evidence to coding agents.  The extrapolation is
plausible because repository-level agent benchmarks show that agents face
cross-file context, integration, and test-feedback problems that resemble human
maintenance work \cite{jimenez-swebench,zhang-codeagent}.  But it remains an
open empirical question whether agents need the same thresholds, respond to the
same smells, or benefit from the same ordering of feedback as human programmers.
The AGENTS.md study is a useful warning: additional repository guidance can
increase cost and sometimes reduce agent success, so agent-facing feedback should
be validated directly rather than inherited uncritically from human practice
\cite{gloaguen-agents}.

The static test-name reference gate is one example of a broader tradeoff:
\texttt{kiss} favors fast proxies that can run in the edit loop over stronger
signals that require project-specific setup.  Runtime coverage, for example,
requires building, configuring, and executing each project's tests, while a
static scan can run across Python and Rust repositories.  The working hypothesis
is that these cheap proxies are better than no structural feedback for
agent-driven refactoring, but that hypothesis needs direct validation on agent
edits rather than support by analogy alone.

\begin{thebibliography}{99}

\bibitem{antinyan}
Vard Antinyan, Miroslaw Staron, and Anna Sandberg.
\newblock Evaluating code complexity triggers, use of complexity measures and
the influence of code complexity on maintenance time.
\newblock \emph{Empirical Software Engineering}, 22:3057--3087, 2017.

\bibitem{elemam}
Khaled El Emam et al.
\newblock The theory of relative dependency.
\newblock \emph{IEEE Software}, 2010.

\bibitem{feathers}
Michael C. Feathers.
\newblock \emph{Working Effectively with Legacy Code}.
\newblock Prentice Hall, 2004.

\bibitem{fowler-code-smell}
Martin Fowler.
\newblock Code Smell.
\newblock \url{https://martinfowler.com/bliki/CodeSmell.html}.

\bibitem{fowler-refactoring}
Martin Fowler.
\newblock \emph{Refactoring: Improving the Design of Existing Code}.
\newblock Addison-Wesley, 2nd edition, 2018.

\bibitem{fowler-quality}
Martin Fowler.
\newblock Is High Quality Software Worth the Cost?
\newblock \url{https://martinfowler.com/articles/is-quality-worth-cost.html}.

\bibitem{fowler-beck}
Martin Fowler.
\newblock Beck Design Rules.
\newblock \url{https://martinfowler.com/bliki/BeckDesignRules.html}.

\bibitem{gerosa}
Marco A. Gerosa et al.
\newblock How do structural dependencies influence change propagation? An
empirical study.
\newblock \emph{International Conference on Software Maintenance}, 2009.

\bibitem{ghaleb}
Tarek A. Ghaleb, Daniel A. da Costa, and Ying Zou.
\newblock An empirical study of the long duration of continuous integration
builds.
\newblock \emph{Empirical Software Engineering}, 24:2102--2139, 2019.

\bibitem{gloaguen-agents}
Gloaguen et al.
\newblock Evaluating AGENTS.md: Are Repository-Level Context Files Helpful for
Coding Agents?
\newblock arXiv:2602.11988, 2026.

\bibitem{inozemtseva}
Laura Inozemtseva and Reid Holmes.
\newblock Coverage is not strongly correlated with test suite effectiveness.
\newblock \emph{International Conference on Software Engineering}, 2014.

\bibitem{jimenez-swebench}
Carlos E. Jimenez et al.
\newblock SWE-bench: Can Language Models Resolve Real-World GitHub Issues?
\newblock \emph{International Conference on Learning Representations}, 2024.

\bibitem{kochhar}
Pavneet Singh Kochhar et al.
\newblock Code coverage and post-release defects: A large-scale study on open
source projects.
\newblock \emph{Mining Software Repositories}, 2017.

\bibitem{mcconnell}
Steve McConnell.
\newblock \emph{Code Complete}.
\newblock Microsoft Press, 2nd edition, 2004.

\bibitem{motwani}
Manas Motwani et al.
\newblock A tale from the trenches: cognitive usefulness and true positives in
static analysis.
\newblock \emph{ICSE Software Engineering in Practice}, 2023.

\bibitem{ousterhout}
John Ousterhout.
\newblock \emph{A Philosophy of Software Design}.
\newblock Yaknyam Press, 2nd edition, 2021.

\bibitem{parnas}
David L. Parnas.
\newblock On the criteria to be used in decomposing systems into modules.
\newblock \emph{Communications of the ACM}, 15(12):1053--1058, 1972.

\bibitem{pragmatic-dry}
Andrew Hunt and David Thomas.
\newblock \emph{The Pragmatic Programmer}.
\newblock Addison-Wesley, 20th anniversary edition, 2019.

\bibitem{rahman-clones}
Foyzur Rahman, Christian Bird, and Premkumar Devanbu.
\newblock Clones: what is that smell?
\newblock \emph{Empirical Software Engineering}, 17:503--530, 2012.

\bibitem{sjoberg}
Dag I. K. Sjoberg, Bente Anda, and Audris Mockus.
\newblock Questioning software maintenance metrics: A comparative case study.
\newblock \emph{Empirical Software Engineering and Measurement}, 2012.

\bibitem{vassallo}
Carmine Vassallo et al.
\newblock How developers engage with static analysis tools in different
contexts.
\newblock \emph{Empirical Software Engineering}, 25:1419--1457, 2020.

\bibitem{zhang-codeagent}
Zhang et al.
\newblock CodeAgent: Enhancing Code Generation with Tool-Integrated Agent
Systems for Real-World Repo-level Coding Challenges.
\newblock \emph{Association for Computational Linguistics}, 2024.

\end{thebibliography}

\end{document}
